\section{Mexanik harakat}

\begin{enumerate}[label=\textbf{\arabic*.}]

\item Mexanik harakat deb nimaga aytiladi?

\begin{enumerate}[label=\Alph*)]
    \item Jism massasining o‘zgarishiga
    \item Jismning boshqa jismlarga nisbatan o‘rnining vaqt davomida o‘zgarishiga
    \item Jismning rang o‘zgarishiga
    \item Jismning hajm o‘zgarishiga
\end{enumerate}



\item Trayektoriya nima?

\begin{enumerate}[label=\Alph*)]
    \item Jismning massasi
    \item Jismning tezligi
    \item Jism harakat davomida chizgan izi
    \item Jismning hajmi
\end{enumerate}



\item Tezlik formulasi qaysi javobda to‘g‘ri berilgan?

\begin{enumerate}[label=\Alph*)]
    \item $v = s \cdot t$
    \item $v = \dfrac{t}{s}$
    \item $v = \dfrac{s}{t}$
    \item $v = s + t$
\end{enumerate}



\item Xalqaro birliklar sistemasida tezlik birligi:

\begin{enumerate}[label=\Alph*)]
    \item m
    \item s
    \item $\text{m/s}$
    \item N
\end{enumerate}



\item Jism 120 m masofani 20 s da bosib o‘tdi. Uning tezligini toping.

\begin{enumerate}[label=\Alph*)]
    \item $4 \ \text{m/s}$
    \item $5 \ \text{m/s}$
    \item $6 \ \text{m/s}$
    \item $8 \ \text{m/s}$
\end{enumerate}



\item Tekis harakat deb qanday harakatga aytiladi?

\begin{enumerate}[label=\Alph*)]
    \item Tezligi o‘zgaradigan harakatga
    \item Trayektoriyasi egri chiziqli bo‘lgan harakatga
    \item Teng vaqt oralig‘ida teng masofa bosib o‘tiladigan harakatga
    \item Faqat vertikal harakatga
\end{enumerate}



\item Jismning bosib o‘tgan yo‘li qanday harf bilan belgilanadi?

\begin{enumerate}[label=\Alph*)]
    \item $t$
    \item $m$
    \item $s$
    \item $v$
\end{enumerate}



\item Vaqtning SI dagi birligi:

\begin{enumerate}[label=\Alph*)]
    \item minut
    \item soat
    \item sekund
    \item sutka
\end{enumerate}



\item Agar jism tezligi ortib borsa, bu qanday harakat deyiladi?

\begin{enumerate}[label=\Alph*)]
    \item Tekis harakat
    \item Notekis harakat
    \item Tinch holat
    \item Aylana harakat
\end{enumerate}



\item Yo‘l formulasi qaysi javobda to‘g‘ri ko‘rsatilgan?

\begin{enumerate}[label=\Alph*)]
    \item $s = vt$
    \item $s = \dfrac{v}{t}$
    \item $s = \dfrac{t}{v}$
    \item $s = v + t$
\end{enumerate}

\end{enumerate}



\textbf{Javoblar:}

\begin{enumerate}
    \item b
    \item c
    \item c
    \item c
    \item c
    \item c
    \item c
    \item c
    \item b
    \item a
\end{enumerate}

